\section{Conclusion}

We presented \method (Adaptive Network of Thought), a split-context prompting method for reasoning over structured, multi-source data.
By decomposing queries into isolated evidence-specific steps and aggregating results with three-valued logic, \method achieves both high accuracy on complex Boolean conditions and robustness to adversarial manipulation---through architectural isolation rather than defense prompts.
On our Yelp-based benchmark, \method outperforms 13 baseline methods by +16\% on clean data while maintaining consistent performance under 8 attack types.

For practitioners deploying LLMs on user-generated content, our key takeaway is that \emph{how} information reaches the model matters as much as what prompting strategy is used.
Split-context processing eliminates entire classes of vulnerabilities by preventing adversarial content from polluting unrelated reasoning steps.

Future work includes extending \method to continuous scoring tasks, exploring hierarchical evidence decomposition for deeply nested queries, and investigating adaptive caching strategies to reduce latency.
We release our benchmark and code to support further research on robust LLM reasoning over multi-source data.
